% GENERATED FILE. Edit canonical lessons, then run npm run generate:books.
% canonical-source-hash: fnv1a64:c5d6c66801e4415e
% canonical-lessons: BN-C475-khoto, BN-C475-gha, BN-C475-mochkano, BN-C475-phola, BN-C475-hochot-khaoya

\chapter{Small Injuries}
\label{ch:bn-small-injuries}
\hlchaptermodality{475}

\begin{chapteropening}
\textbf{By the end of this chapter:} \emph{I can talk about wounds, sores, sprains and stumbling.}

Complete the last lesson of chapter 475: I can talk about wounds, sores, sprains and stumbling.
\end{chapteropening}

\section[\texorpdfstring{\bn{ক্ষত} (khôtô)}{ক্ষত (khôtô)}]{\bn{ক্ষত} (khôtô) — a wound}

\label{lesson:BN-C475-khoto}



\begin{quote}
Before the new one: say the Bengali for cancer, then the Bengali for jaundice.
\end{quote}

\subsection*{You'll want to know: \bn{ক্ষত}}
\textbf{\bn{ক্ষত}} — \emph{khôtô} — \textquotedblleft{}a wound\textquotedblright{}.

\begin{quote}
\bn{নার্স} \bn{ক্ষতটায়} \bn{মলম} \bn{লাগালেন।}
\end{quote}

\begin{quote}
\emph{nārs khôtôṭāy môlôm lāgālen.}
\end{quote}

\begin{quote}
The nurse put ointment on the wound.
\end{quote}

\begin{cousinweb}
\textbf{\bn{ক্ষত}} (\emph{khôtô}) is a tatsama, Sanskrit \emph{kṣata}, \textquotedblleft{}wounded\textquotedblright{}, and as a noun \textquotedblleft{}a wound\textquotedblright{} (CDIAL 3645).
\end{cousinweb}

\begin{grammarlens}[title={What you've built}]
One more word for describing what you lived and felt.
\end{grammarlens}

\subsection*{Your turn}
\begin{itemize}
\raggedright
  \item \emph{Say it:} \emph{khôtô}
  \item \emph{Say it:} \emph{khôtô}, once more
  \item \emph{Say it:} \emph{nārs khôtôṭāy môlôm lāgālen.} — The nurse put ointment on the wound.
  \item \emph{Recall:} say the Bengali for cancer, then the Bengali for jaundice, then say \emph{khôtô} again
\end{itemize}

\begin{tcolorbox}[breakable,colback=teal!4,colframe=teal!35!black,title={Before you move on}]
What does \bn{ক্ষত} mean? (\textquotedblleft{}A wound\textquotedblright{}.) Say it once more.
\end{tcolorbox}
\section[\texorpdfstring{\bn{ঘা} (ghā)}{ঘা (ghā)}]{\bn{ঘা} (ghā) — a sore; an ulcer}

\label{lesson:BN-C475-gha}



\begin{quote}
Before the new one: say the Bengali for jaundice, then the Bengali for a wound.
\end{quote}

\subsection*{You'll want to know: \bn{ঘা}}
\textbf{\bn{ঘা}} — \emph{ghā} — \textquotedblleft{}a sore; an ulcer\textquotedblright{}.

\begin{quote}
\bn{অরুণের} \bn{হাতে} \bn{একটা} \bn{ঘা} \bn{হল।}
\end{quote}

\begin{quote}
\emph{ôruner hāte ekṭā ghā hôlô.}
\end{quote}

\begin{quote}
Arun got a sore on his hand.
\end{quote}

\begin{cousinweb}
\textbf{\bn{ঘা}} (\emph{ghā}) is listed by CDIAL for Bengali, as a wound or sore, under Sanskrit \emph{ghāta}, \textquotedblleft{}a blow\textquotedblright{} (CDIAL 4460).
\end{cousinweb}

\begin{grammarlens}[title={What you've built}]
One more word for describing what you lived and felt.
\end{grammarlens}

\subsection*{Your turn}
\begin{itemize}
\raggedright
  \item \emph{Say it:} \emph{ghā}
  \item \emph{Say it:} \emph{ghā}, once more
  \item \emph{Say it:} \emph{ôruner hāte ekṭā ghā hôlô.} — Arun got a sore on his hand.
  \item \emph{Recall:} say the Bengali for jaundice, then the Bengali for a wound, then say \emph{ghā} again
\end{itemize}

\begin{tcolorbox}[breakable,colback=teal!4,colframe=teal!35!black,title={Before you move on}]
What does \bn{ঘা} mean? (\textquotedblleft{}A sore; an ulcer\textquotedblright{}.) Say it once more.
\end{tcolorbox}
\section[\texorpdfstring{\bn{মচকানো} (môchkāno)}{মচকানো (môchkāno)}]{\bn{মচকানো} (môchkāno) — to sprain; to twist (a joint)}

\label{lesson:BN-C475-mochkano}



\begin{quote}
Before the new one: say the Bengali for a wound, then the Bengali for a sore.
\end{quote}

\subsection*{You'll want to know: \bn{মচকানো}}
\textbf{\bn{মচকানো}} — \emph{môchkāno} — \textquotedblleft{}to sprain; to twist (a joint)\textquotedblright{}.

\begin{quote}
\bn{সিঁড়িতে} \bn{মীরার} \bn{গোড়ালি} \bn{মচকাল।}
\end{quote}

\begin{quote}
\emph{sĩṛite mīrār goṛāli môchkālô.}
\end{quote}

\begin{quote}
Mira twisted her ankle on the stairs.
\end{quote}

Past (he/she …): \bn{সেই} \bn{খেলোয়াড়ের} \bn{হাঁটু} \bn{মচকাল}, \emph{sei kheloyāṛer hā̃ṭu môchkālô}.

\begin{cousinweb}
\textbf{\bn{মচকানো}} (\emph{môchkāno}) is a Bengali verb for twisting a joint whose origin is not settled.
\end{cousinweb}

\begin{grammarlens}[title={What you've built}]
One more word for describing what you lived and felt.
\end{grammarlens}

\subsection*{Your turn}
\begin{itemize}
\raggedright
  \item \emph{Say it:} \emph{môchkāno}
  \item \emph{Say it:} \emph{môchkāno}, once more
  \item \emph{Say it:} \emph{sĩṛite mīrār goṛāli môchkālô.} — Mira twisted her ankle on the stairs.
  \item \emph{Recall:} say the Bengali for a wound, then the Bengali for a sore, then say \emph{môchkāno} again
\end{itemize}

\begin{tcolorbox}[breakable,colback=teal!4,colframe=teal!35!black,title={Before you move on}]
What does \bn{মচকানো} mean? (\textquotedblleft{}To sprain; to twist (a joint)\textquotedblright{}.) Say it once more.
\end{tcolorbox}
\section[\texorpdfstring{\bn{ফোলা} (pholā)}{ফোলা (pholā)}]{\bn{ফোলা} (pholā) — to swell}

\label{lesson:BN-C475-phola}



\begin{quote}
Before the new one: say the Bengali for a sore, then the Bengali for to sprain.
\end{quote}

\subsection*{You'll want to know: \bn{ফোলা}}
\textbf{\bn{ফোলা}} — \emph{pholā} — \textquotedblleft{}to swell\textquotedblright{}.

\begin{quote}
\bn{মচকানোর} \bn{পরে} \bn{গোড়ালিটা} \bn{ফুলে} \bn{গেল।}
\end{quote}

\begin{quote}
\emph{môchkānor pôre goṛāliṭā phule gelô.}
\end{quote}

\begin{quote}
After the sprain the ankle swelled up.
\end{quote}

Past (he/she …): \bn{সেই} \bn{রোগীর} \bn{গাল} \bn{ফুলল}, \emph{sei rogīr gāl phullô}.

\begin{cousinweb}
\textbf{\bn{ফোলা}} (\emph{pholā}) is listed by CDIAL for Bengali, as swelling and blossoming, under Sanskrit \emph{phullati}, \textquotedblleft{}opens out like a flower\textquotedblright{} (CDIAL 9093).
\end{cousinweb}

\begin{grammarlens}[title={What you've built}]
One more word for describing what you lived and felt.
\end{grammarlens}

\subsection*{Your turn}
\begin{itemize}
\raggedright
  \item \emph{Say it:} \emph{pholā}
  \item \emph{Say it:} \emph{pholā}, once more
  \item \emph{Say it:} \emph{môchkānor pôre goṛāliṭā phule gelô.} — After the sprain the ankle swelled up.
  \item \emph{Recall:} say the Bengali for a sore, then the Bengali for to sprain, then say \emph{pholā} again
\end{itemize}

\begin{tcolorbox}[breakable,colback=teal!4,colframe=teal!35!black,title={Before you move on}]
What does \bn{ফোলা} mean? (\textquotedblleft{}To swell\textquotedblright{}.) Say it once more.
\end{tcolorbox}
\section[\texorpdfstring{\bn{হোঁচট} \bn{খাওয়া} (hõchôṭ khāoyā)}{হোঁচট খাওয়া (hõchôṭ khāoyā)}]{\bn{হোঁচট} \bn{খাওয়া} (hõchôṭ khāoyā) — to stumble; to trip}

\label{lesson:BN-C475-hochot-khaoya}



\begin{quote}
Before the new one: say the Bengali for to sprain, then the Bengali for to swell.
\end{quote}

\subsection*{You'll want to know: \bn{হোঁচট} \bn{খাওয়া}}
\textbf{\bn{হোঁচট} \bn{খাওয়া}} — \emph{hõchôṭ khāoyā} — \textquotedblleft{}to stumble; to trip\textquotedblright{}.

\begin{quote}
\bn{রাতে} \bn{রাস্তায়} \bn{পাথরে} \bn{হোঁচট} \bn{খেয়ে} \bn{অরুণ} \bn{পড়ে} \bn{গেল।}
\end{quote}

\begin{quote}
\emph{rāte rāstāy pāthôre hõchôṭ kheye ôrun pôṛe gelô.}
\end{quote}

\begin{quote}
At night Arun tripped on a stone in the road and fell.
\end{quote}

Past (he/she …): \bn{তিনি} \bn{সিঁড়িতে} \bn{হোঁচট} \bn{খেলেন}, \emph{tini sĩṛite hõchôṭ khelen}.

\begin{cousinweb}
\textbf{\bn{হোঁচট} \bn{খাওয়া}} (\emph{hõchôṭ khāoyā}) is \emph{hõchôṭ}, \textquotedblleft{}a stumble\textquotedblright{}, whose origin is not settled, with \emph{khāoyā}, \textquotedblleft{}to eat\textquotedblright{} (CDIAL 3865): Bengali \textquotedblleft{}eats\textquotedblright{} a stumble as it eats a scolding.
\end{cousinweb}

\begin{grammarlens}[title={What you've built}]
One more word for describing what you lived and felt.
\end{grammarlens}

\subsection*{Your turn}
\begin{itemize}
\raggedright
  \item \emph{Say it:} \emph{hõchôṭ khāoyā}
  \item \emph{Say it:} \emph{hõchôṭ khāoyā}, once more
  \item \emph{Say it:} \emph{rāte rāstāy pāthôre hõchôṭ kheye ôrun pôṛe gelô.} — At night Arun tripped on a stone in the road and fell.
  \item \emph{Recall:} say the Bengali for to sprain, then the Bengali for to swell, then say \emph{hõchôṭ khāoyā} again
\end{itemize}

\begin{tcolorbox}[breakable,colback=teal!4,colframe=teal!35!black,title={Before you move on}]
What does \bn{হোঁচট} \bn{খাওয়া} mean? (\textquotedblleft{}To stumble; to trip\textquotedblright{}.) Say it once more.
\end{tcolorbox}
